% Bagian: first-order-logic
% Bab: introduction
% Subbagian: first-order-logic

\documentclass[../../../include/open-logic-section]{subfiles}

\begin{document}

\olfileid[id]{fol}{int}{fol}

\olsection{Logika Orde Pertama}

Anda mungkin sudah mengenal logika orde pertama dari pengantar pertama Anda
tentang logika formal.\footnote{Bahkan, kurang lebih kami menganggap Anda
sudah mengenalnya!{} Jika belum, Anda dapat meninjau buku teks yang lebih
elementer, seperti \emph{forall x} \citep{Magnus2021}.} Anda mungkin
mengenalnya sebagai ``logika kuantifikasional'' atau ``logika predikat.''
Pertama-tama, logika orde pertama merupakan suatu bahasa formal. Artinya,
bahasa ini mempunyai kosakata tertentu, dan ungkapan-ungkapannya merupakan
untai yang tersusun dari kosakata tersebut. Namun, tidak setiap untai
diperbolehkan. Ada beberapa jenis ungkapan yang diperbolehkan: suku,
!!{formula}s, dan !!{sentence}s. Kita terutama berminat pada !!{sentence}s
logika orde pertama: kalimat-kalimat tersebut memberi kita padanan formal
bagi kalimat bahasa Inggris, dan mengenainya kita dapat mengajukan
pertanyaan-pertanyaan yang lazim diminati ahli logika. Misalnya:
\begin{itemize}
    \item Apakah $!B$ mengikuti secara logis dari~$!A$?
    \item Apakah $!A$ benar secara logis, salah secara logis, atau
kontingen?
    \item Apakah $!A$ dan $!B$ ekuivalen?
\end{itemize}

Pertanyaan-pertanyaan ini terutama menyangkut ``makna'' !!{sentence}s logika
orde pertama. Misalnya, seorang filsuf akan menganalisis pertanyaan apakah
$!B$ mengikuti secara logis dari~$!A$ dengan bertanya: adakah suatu kasus
ketika $!A$ benar tetapi~$!B$ salah ($!B$ tidak mengikuti dari~$!A$), atau
apakah setiap kasus yang membuat $!A$ benar juga membuat~$!B$ benar ($!B$
memang mengikuti dari~$!A$)? Namun, kita belum diberi tahu apa yang dimaksud
dengan suatu ``kasus''---itulah tugas \emph{semantik}. Semantik logika orde
pertama menyediakan model yang tepat secara matematis bagi gagasan intuitif
filsuf tentang ``kasus,'' dan juga---ini penting---bagi apa artinya suatu
!!a{sentence}~$!A$ \emph{benar dalam} suatu kasus. Model matematis yang tepat
yang akan kita kembangkan disebut !!a{structure}. Relasi yang memperjelas
makna ``benar dalam'' disebut relasi \emph{keterpenuhan}. Jadi, yang akan kita
definisikan ialah ``$!A$ terpenuhi dalam~$\Struct{M}$'' (dalam simbol:
$\Sat{M}{!A}$) untuk !!{sentence}s~$!A$ dan !!{structure}s~$\Struct{M}$.
Setelah ini dilakukan, kita juga dapat memberikan definisi yang tepat bagi
istilah-istilah semantis lain seperti ``mengikuti dari'' atau ``benar secara
logis.'' Definisi-definisi tersebut memungkinkan kita menentukan, sekali
lagi dengan ketepatan matematis, apakah, misalnya,
$\lforall[x][(!A(x) \lif !B(x))],
\lexists[x][!A(x)] \Entails \lexists[x][!B(x)]$. Tentu saja jawabannya
``ya.'' Jika Anda sudah dilatih menyimbolkan kalimat bahasa Inggris dalam
logika orde pertama, Anda akan mengenalinya sebagai, misalnya, penyimbolan
``Semua semut adalah serangga; ada semut; maka ada serangga.'' Argumen itu
jelas valid, sehingga model matematis kita bagi ``mengikuti dari'' dalam
bahasa formal harus memberikan jawaban yang sama.

Topik lain yang mungkin Anda ingat dari pengantar pertama tentang logika
formal ialah adanya \emph{!!{derivation}s}. Jika Anda pernah mengikuti mata
kuliah pengantar logika formal, pengajar Anda mungkin meminta Anda berlatih
menemukan !!{derivation}s semacam itu, bahkan mungkin !!a{derivation} yang
menunjukkan bahwa perikutan di atas berlaku. Ada banyak cara berbeda untuk
memberikan !!{derivation}s: Anda mungkin pernah mengerjakan sesuatu yang
disebut ``deduksi alami'' atau ``pohon kebenaran,'' tetapi masih banyak cara
lain. Tujuan sistem !!{derivation} ialah menyediakan alat yang dapat digunakan
untuk menjawab pertanyaan-pertanyaan ahli logika di atas: misalnya, suatu
!!{derivation} deduksi alami dengan
$\lforall[x][(!A(x) \lif !B(x))]$ dan $\lexists[x][!A(x)]$ sebagai premis,
serta $\lexists[x][!B(x)]$ sebagai kesimpulan (baris terakhir),
\emph{memverifikasi} bahwa $\lexists[x][!B(x)]$ mengikuti secara logis dari
$\lforall[x][(!A(x) \lif !B(x))]$ dan $\lexists[x][!A(x)]$.

Namun, mengapa demikian? Sepintas lalu, sistem !!{derivation} sama sekali
tidak berkaitan dengan semantik: memberikan suatu !!{derivation} formal hanya
melibatkan penyusunan simbol dengan cara tertentu yang diatur oleh aturan;
sistem itu tidak menyebut ``kasus'' atau ``benar dalam'' sama sekali. Hubungan
antara sistem !!{derivation} dan semantik harus ditetapkan melalui suatu
penyelidikan metalogis. Yang diperlukan ialah bukti matematis, misalnya,
bahwa suatu !!{derivation} formal atas $\lexists[x][!B(x)]$ dari premis
$\lforall[x][(!A(x) \lif !B(x))]$ dan $\lexists[x][!A(x)]$ dimungkinkan jika,
dan hanya jika, $\lforall[x][(!A(x) \lif !B(x))]$ dan
$\lexists[x][!A(x)]$ bersama-sama mengakibatkan~$\lexists[x][!B(x)]$.
Namun, sebelum hal ini dapat dilakukan, banyak pekerjaan yang teliti harus
dilaksanakan agar definisi sintaksis dan semantik benar.

\end{document}
